\subsection{子模的直和}

定义 6. 称 A-模 E 是它的子模的一个族 \((\mathbf{M}_{\iota})_{\iota \in I}\) 的直和，如果典范映射 \(\bigoplus_{\iota \in I} \mathbf{M}_{\iota} \to \mathbf{E}\) （第 6 号）是同构。

这等价于说，每个 \(x \in \mathbf{E}\) 都可以以唯一的方式写成形式 \(x = \sum_{i \in I} x_i\) ，其中 \(x_i \in \mathbf{E}_i\) 对所有 \(i \in I\) 成立；元素 \(x_i\) （这样与 \(x\) 对应的元素）称为 \(x\) 在 \(\mathbf{E}_i\) 中的分量；映射 \(x \mapsto x_i\) 是线性的。

注记。(1) 设 \((\mathbf{M}_{\iota})_{\iota \in \mathbf{I}}, (\mathbf{N}_{\iota})_{\iota \in \mathbf{I}}\) 是模 \(\mathbf{E}\) 的两个子模族，具有同一指标集；假设 \(\mathbf{E}\) 既是族 \((\mathbf{M}_{\iota})\) 的直和又是族 \((\mathbf{N}_{\iota})\) 的直和，且 \(\mathbf{N}_{\iota} \subset \mathbf{M}_{\iota}\) （对所有 \(\iota \in \mathbf{I}\) ）。则 \(\mathbf{N}_{\iota} = \mathbf{M}_{\iota}\) 对所有 \(\iota \in \mathbf{I}\) 成立，这由第 6 号命题 7 的推论 1 应用于典范单射 \(f_{\iota}: \mathbf{N}_{\iota} \to \mathbf{M}_{\iota}\) 立即得出。

命题 11. 设 \((\mathbf{M}_{i})_{i\in I}\) 是 A-模 E 的一族子模。下列性质等价：

\begin{enumerate}
\def\labelenumi{(\alph{enumi})}
\item
  子模 \(\sum_{i\in I}\mathbf{M}_i\) 是族 \((\mathbf{M}_i)_{i\in I}\) 的直和。
\item
  关系 \(\sum_{i\in I}x_i = 0\) （其中 \(x_{i}\in \mathbf{M}_{i}\) 对所有 \(i\in I\) 成立）蕴含 \(x_{i} = 0\) 对所有 \(i\in I\) 成立
\item
  对所有 \(\kappa \in \mathbf{I}\) ， \(\mathbf{M}_{\kappa}\) 与 \(\sum_{i \neq \kappa} \mathbf{M}_i\) 的交约化为 0。
\end{enumerate}

(a) 与 (b) 等价是直接的，因为关系 \(\sum_{i} x_{i} = \sum_{i} y_{i}\) 等价于 \(\sum_{i} (x_{i} - y_{i}) = 0\) 。另一方面，根据定义 6，由 \(\bigoplus_{i \in I} M_{i}\) 的元素作为元素 \(x_{i} \in M_{i}\) 的直和的表达式的唯一性，(a) 蕴含 (c)。最后，关系 \(\sum_{i} x_{i} = 0\) （其中 \(x_{i} \in M_{i}\) 对所有 \(i\) 成立）可以写成：对所有 \(\kappa \in I\) ， \(x_{\kappa} = \sum_{i \neq \kappa} (-x_{i})\) ；于是条件 (c) 蕴含 \(x_{\kappa} = 0\) 对所有 \(\kappa \in I\) 成立，从而 (c) 蕴含 (b)。

定义 7. 自同态 \(e\) （A-模 \(\mathbf{E}\) 上的自同态）称为投影算子，如果 \(e \circ e = e\) （换言之，如果 \(e\) 是环 \(\operatorname{End}(\mathbf{E})\) 中的幂等元）。在 \(\operatorname{End}(\mathbf{E})\) 中，投影算子的一族 \((e_{\lambda})_{\lambda \in L}\) 称为正交的，如果 \(e_{\lambda} \circ e_{\mu} = 0\) （对 \(\lambda \neq \mu\) ）。

命题 12. 设 E 是一个 A-模。

\begin{enumerate}
\def\labelenumi{(\roman{enumi})}
\item
  若 \(\mathbf{E}\) 是子模族 \((\mathbf{M}_{\lambda})_{\lambda \in \mathbf{L}}\) 的直和，且对所有 \(x \in \mathbf{E}\) ， \(e_{\lambda}(x)\) 是 \(x\) 在 \(\mathbf{M}_{\lambda}\) 中的分量，则 \((e_{\lambda})\) 是一个正交投影算子族，使得 \(x = \sum_{\lambda \in \mathbf{L}} e_{\lambda}(x)\) 对所有 \(x \in \mathbf{E}\) 成立。
\item
  反之，若 \((e_{\lambda})_{\lambda \in \mathbf{L}}\) 是 \(\operatorname{End}(\mathbf{E})\) 中的一个正交投影算子族，使得 \(x = \sum_{\lambda \in \mathbf{L}} e_{\lambda}(x)\) 对所有 \(x \in \mathbf{E}\) 成立，则 \(\mathbf{E}\) 是子模族 \(\mathbf{M}_{\lambda} = e_{\lambda}(\mathbf{E})\) 的直和。
\end{enumerate}

性质 (i) 由定义得出，而 (ii) 是第 6 号命题 6 的推论 2 应用于典范单射 \(\mathbf{M}_{\lambda} \to \mathbf{E}\) 和映射 \(e_{\lambda}: \mathbf{E} \to \mathbf{M}_{\lambda}\) 的一个特殊情形。

注意，当 \(\mathbf{L}\) 有限时，条件 \(x = \sum_{\lambda \in \mathbf{L}} e_{\lambda}(x)\) （对所有 \(x \in \mathbf{E}\) ）在 \(\operatorname{End}(\mathbf{E})\) 中也可以写成

\[
1 _ {\mathbf {E}} = \sum_ {\lambda \in L} e _ {\lambda}.\tag{31}
\]

推论。设 \(e\) 是 \(\mathbf{E}\) 的一个投影算子，则 \(\mathbf{E}\) 是像 \(\mathbf{M} = e(\mathbf{E})\) 与核 \(\mathbf{N} = e^{-1}(0)\) （即 \(e\) 的像与核）的直和；对所有 \(x = x_1 + x_2 \in \mathbf{E}\) （其中 \(x_1 \in \mathbf{M}\) ， \(x_2 \in \mathbf{N}\) ），有 \(x_1 = e(x)\) ；1 - \(e\) 是 \(\mathbf{E}\) 的一个投影算子，其像为 \(\mathbf{N}\) ，核为 \(\mathbf{M}\) 。

\((1 - e)^2 = 1 - 2e + e^2 = 1 - e\) 在 \(\operatorname{End}(\mathbf{E})\) 中成立，因此 \(1 - e\) 是投影算子；又由于 \(e(1 - e) = (1 - e)e = e - e^2 = 0\) ，根据命题 12， \(\mathbf{E}\) 是像 \(\mathbf{M}\) 和 \(\mathbf{N}\) （分别为 \(e\) 与 \(1 - e\) 的像）的直和。最后，对所有 \(x \in \mathbf{E}\) ，关系 \(x \in \mathbf{M}\) 等价于 \(x = e(x)\) ；因为 \(x = e(x)\) 按定义蕴含 \(x \in \mathbf{M}\) ，而反之，若 \(x = e(x')\) （其中 \(x' \in \mathbf{E}\) ），则 \(e(x) = e^2(x') = e(x') = x\) ；

这就表明 M 是 1 - e 的核；交换 e 与 1 - e 的角色，类似地可以看出 N 是 e 的核。

注记。(2) 设 E, F 是两个 A-模，使得 E 是子模的一个有限族 \((\mathbf{M}_{i})_{1\leqslant i\leqslant m}\) 的直和，而 F 是子模的一个有限族 \((\mathbf{N}_{j})_{1\leqslant j\leqslant n}\) 的直和。已知（第 6 号，命题 6 的推论 1） \(\operatorname{Hom}_{\mathbf{A}}(\mathbf{E},\mathbf{F})\) 典范地等同于乘积 \(\prod_{i,j}\operatorname{Hom}_{\mathbf{A}}(\mathbf{M}_{i},\mathbf{N}_{j})\) ；确切地说，对每个族 \((u_{ji})\) （其中 \(u_{ji}\in\operatorname{Hom}_{\mathbf{A}}(\mathbf{M}_{i},\mathbf{N}_{j})\) ），对应着按下述方式定义的线性映射 \(u:E\to F\) 。只需定义 u 在每个 \(M_{i}\) 上的限制，并且对每个 \(x_{i}\in M_{i}\) ，

\[
u (x _ {i}) = \sum_ {j = 1} ^ {n} u _ {j i} (x _ {i}).
\]

现在设 \(\mathbf{G}\) 是第三个 A-模，是子模的一个有限族 \((\mathbf{P}_k)_{1 \leqslant k \leqslant p}\) 的直和；设 \(v\) 是 \(\mathbf{F}\) 到 \(\mathbf{G}\) 中的线性映射，并设 \((v_{kj}) \in \prod_{j,k} \operatorname{Hom}_{\mathbf{A}}(\mathbf{N}_j, \mathbf{P}_k)\) 是典范地对应于它的族。对所有 \(x_i \in \mathbf{M}_i\) ，

\[
v (u (x _ {i})) = \sum_ {j = 1} ^ {n} v (u _ {j i} (x _ {i})) = \sum_ {k = 1} ^ {p} \sum_ {j = 1} ^ {n} v _ {k j} (u _ {j i} (x _ {i})).
\]

由此可见，若我们写

\[
w _ {k i} = \sum_ {j = 1} ^ {n} v _ {k j} \circ u _ {j i} \in \operatorname{Hom} _ {\mathbf {A}} (\mathrm{M} _ {i}, \mathrm{P} _ {k})\tag{32}
\]

族 \((w_{ki})\) 典范地对应于复合线性映射 \(w = v \circ u\) （从 \(\mathbf{E}\) 到 \(\mathbf{G}\) ）（参见 §10，第 5 号）。

